\section{项目背景与目标}

金融市场仿真器是金融研究和量化交易实践的重要工具，它能够在不影响真实市场的情况下模拟交易行为、市场机制以及政策干预的影响，为策略验证、风险评估和政策制定提供可控的实验环境。传统的研究方法包括实盘检验和实验室小市场试验，但受限于数据隐私、实验成本和可控性，这些方法往往难以观察个体交易者的行为或进行大规模的“假设分析”。计算机仿真弥补了这一不足，使研究者能够完全控制市场规则、观察每个代理的动作，并且重复实验以评估不同假设的影响。

ABIDES（Agent‐Based Interactive Discrete Event Simulation）是当前开源社区中较为完善的高保真市场仿真平台。它采用离散事件驱动内核，支持纳秒级时间分辨率、消息传递式通信、可配置网络延迟以及基于 ITCH/OUCH 协议的交易所接口。ABIDES 在模拟单一股票的连续双向拍卖及多智能体交互方面表现优秀，并提供了丰富的基础代理类和示例策略。然而，ABIDES 仍存在以下局限：
\begin{itemize}
  \item \textbf{单资产限制：}ABIDES 的交易所和订单簿按单个股票设计，缺乏跨股票或跨资产的统一市场架构，使得无法研究资产间的联动效应。
  \item \textbf{缺乏跨日机制：}现有系统主要支持日内连续竞价交易，尚未实现开盘集合竞价等跨日场景。
  \item \textbf{依赖外部校准：}背景代理的行为参数需要人工或外部算法反复调节，缺乏内置的自校准机制，难以在保持高效的同时逼近真实市场统计特征。
  \item \textbf{性能瓶颈：}由于采用单线程 Python 实现，尽管可以模拟上万次事件，但在多标的和海量代理场景下难以满足全市场仿真的吞吐需求。
\end{itemize}

为了克服上述不足，本项目计划在 ABIDES 平台的基础上研发新一代全尺寸金融市场仿真器，实现以下目标：
\begin{enumerate}
  \item \textbf{跨资产交互：}支持多个交易品种和交易所的统一仿真，允许代理同时在多只股票或品类上发单，模拟现实中的跨资产套利、行业传播等现象。
  \item \textbf{跨日模拟：}引入连续竞价与集合竞价的切换机制，支持多日交易的状态传递、隔夜持仓以及开盘价形成过程。
  \item \textbf{动态自校准：}设计内置的自校准模块，使背景代理能够根据历史数据或实时统计指标自动调整行为，无需大量外部迭代即可逼近真实市场的波动和流动性特征。
  \item \textbf{高效可扩展：}改进订单簿数据结构和撮合算法，引入多线程或异步架构，提升在千股级别和万级代理下的仿真吞吐量。
  \item \textbf{易用性与可视化：}提供灵活的配置系统和交互接口，支持用户实时插入策略、修改政策参数，并通过可视化界面监控市场动态和代理表现。
\end{enumerate}

通过实现以上目标，新仿真器将为研究者和从业者提供一个可用于大规模、跨资产、高保真市场模拟的平台，便于开展 AI 交易研究、策略评估、政策测试以及数据生成等多方面工作。